%% Global Plastic Pollution & Ocean Health Dashboard --- Final Report
\documentclass[10pt,twocolumn]{article}

\usepackage[top=1.8cm,bottom=1.8cm,left=1.6cm,right=1.6cm,columnsep=0.6cm]{geometry}
\usepackage[T1]{fontenc}
\usepackage[utf8]{inputenc}
\usepackage{amsmath}
\usepackage{booktabs}
\usepackage{tabularx}
\usepackage{xcolor}
\usepackage{url}
\usepackage{graphicx}

\setlength{\parskip}{3pt}
\setlength{\parindent}{0pt}
\setlength{\abovecaptionskip}{4pt}
\setlength{\belowcaptionskip}{2pt}

\newenvironment{tightitem}{%
  \begin{itemize}%
    \setlength{\itemsep}{2pt}%
    \setlength{\parskip}{0pt}%
    \setlength{\parsep}{0pt}%
    \setlength{\topsep}{3pt}%
}{\end{itemize}}

\newcommand{\act}[1]{\textbf{Act\,#1}}
\newcommand{\tool}[1]{\texttt{#1}}

% ─────────────────────────────────────────────────────────────────────────────
\begin{document}

\twocolumn[{%
\begin{center}
  {\large\bfseries
    Global Plastic Pollution and Ocean Health Dashboard\\[4pt]}
  {\normalsize A Three-Act Interactive Narrative Visualization}\\[5pt]
  {\small Nguyen Huy Hung \quad \texttt{25hung.nh2@vinuni.edu.vn} \quad
   Nguyen Ba Thanh Bac \quad \texttt{25bac.nbt@vinuni.edu.vn}}\\[2pt]
  {\small VinUniversity}\\[3pt]
  {\small June 2026}
\end{center}
\smallskip
\noindent\rule{\linewidth}{0.5pt}
\begin{center}
\begin{minipage}{0.93\linewidth}
  \small\textbf{Abstract.}
  We present a full-stack interactive dashboard that narrates the journey of
  plastic pollution from land-based sources to ocean degradation across three
  linked analytical acts.
  The first act visualises mismanaged plastic waste per capita for 200+
  countries from 1990 to 2019.
  The second reveals 1\,000+ river emission points spatially clustered via
  emission-weighted K-Means, exposing the geographic concentration of
  riverine plastic loading.
  The third act explores the Ocean Health Index (OHI) for 220+ coastal
  regions from 2012 to 2023, equipped with four statistical and machine-learning
  forecast models.
  The system provides a React + D3 production frontend, a Python/FastAPI
  REST backend, and an equivalent Python Shiny research application, sharing
  the same preprocessing and ML pipeline.
\end{minipage}
\end{center}
\noindent\rule{\linewidth}{0.5pt}
\vspace{6pt}
}]

%% ────────────────────────────────────────────────────────────────────────────
\section{Motivation and Problem Statement}

Approximately eight million metric tonnes of plastic enter the world's oceans
every year [1].
This contamination disrupts marine food webs, degrades coastal ecosystems,
and imposes economic costs estimated in the billions of dollars annually.
Despite growing public awareness, most data-driven tools treat plastic
pollution as an isolated statistic: they show \emph{where} mismanagement is
worst, but do not trace the causal pathway from land-based waste through
riverine transport to actual ocean health degradation.

The goal of this project is to build a \emph{narrative data visualization}
that connects these three stages in a single, coherent analytical journey.
Following Segel and Heer's taxonomy of narrative structures [5], we adopt
an \emph{author-driven} macro-layout that imposes a deliberate sequence
(country $\to$ rivers $\to$ ocean) while providing \emph{reader-driven}
micro-interactions within each stage.
This dual structure ensures novice users receive a guided story while
domain experts can freely explore the data.

The intended audience spans environmental policy analysts who need
country-level summaries, hydrologists studying river contributions, and
marine scientists tracking ocean health trends.
Each group is served by one of the three acts, and all benefit from the
continuous visual thread (consistent map projections, shared colour language,
and sticky navigation) that links the acts together.

%% ────────────────────────────────────────────────────────────────────────────
\section{Data Sources and Preprocessing}

\subsection{Dataset Overview}

The project integrates three public datasets summarised in
Table~\ref{tab:datasets}.

\begin{table}[ht]
\centering
\small
\caption{Datasets used across the three acts.}
\label{tab:datasets}
\resizebox{\columnwidth}{!}{%
\begin{tabular}{@{}p{2.1cm}p{1.6cm}rp{1.7cm}@{}}
\toprule
\textbf{Dataset} & \textbf{Source} & \textbf{Records} & \textbf{Scope}\\
\midrule
Mismanaged plastic waste per capita (kg/person/yr)
  & OWID [4] & $\sim$5\,700 & 200+ countries, 1990--2019\\
\addlinespace[2pt]
River plastic emission points (tonnes/yr)
  & Meijer et al.\ [2] & $\sim$32\,000 & 1\,000+ rivers, shapefile\\
\addlinespace[2pt]
Ocean Health Index scores (0--100)
  & OHI [3] & 341\,125 & 220+ regions, 2012--2023\\
\bottomrule
\end{tabular}}
\end{table}

\subsection{Plastic Waste Preprocessing}

The Our World in Data CSV is distributed with both country-level rows and
supra-national aggregate rows (identified by entity codes beginning
\tool{OWID\_}).
These aggregates were filtered out to retain only sovereign-state
observations, giving a clean set of per-country, per-year waste rates.
A unified four-column schema (\tool{entity}, \tool{code}, \tool{year},
\tool{value}) was enforced for all downstream queries.
The field \tool{code} stores ISO\,3166-1 alpha-3 codes, which Plotly can
directly consume for choropleth rendering via
\tool{locationmode='ISO-3'}.

\subsection{River Emissions Preprocessing}

The Meijer dataset is distributed as a zipped ESRI shapefile.
At application startup, the archive is extracted to a temporary directory,
read with \tool{geopandas}, and the process directory is immediately deleted
to avoid disk bloat.
Point coordinates are derived from \tool{geometry.x} and \tool{geometry.y}.
Rivers with zero or negative recorded emission values are dropped
($\approx$0.3\,\% of records).

A spatial join using \tool{geopandas.sjoin} with the \texttt{within} predicate
against Natural Earth 110\,m administrative country boundaries enriches each
river record with country name, ISO\,A3 code, and continent.
Records that fall outside any polygon (open ocean) receive \texttt{country\,=\,"Unknown"} and are excluded from country-level aggregations but retained on the
map.
Duplicate join hits (rivers near borders) are resolved by keeping the first
match, consistent with the original dataset's usage.

\subsection{OHI Score Preprocessing}

The 341\,125-row CSV stores scores across multiple goals
(the composite \emph{Index} plus 10 sub-goals), multiple dimensions
(\emph{score}, \emph{future}, \emph{status}), and every year from 2012 to
2023.
Column names were standardised and both \tool{year} and \tool{region\_id}
were cast to integers.
Region\,0 (\emph{Global average}) is kept in the dataset but masked from the
choropleth map and handled as a special aggregation case in trend and
forecast queries.
A deduplicated goal-metadata table (\tool{goal, long\_goal}) was extracted
at load time to populate the \emph{Goal / Sub-Goal} dropdown without
repeated runtime filtering.

%% ────────────────────────────────────────────────────────────────────────────
\section{Visualization and Interaction Design}

\subsection{Narrative Architecture}

The three-act structure is implemented as a tabbed page navigation (React:
sticky section headers with IntersectionObserver; Shiny: \tool{page\_navbar}).
Within each act, a \emph{sticky sidebar} holds all interactive controls, while
the main panel contains charts that grow to their natural height and scroll
independently.
This layout mirrors best practices in analytical dashboard design: controls
are always visible without interrupting the reading flow of the charts below.

\subsection{Act I --- Country-Level Plastic Crisis}

\textbf{Choropleth world map.}
Mismanaged plastic waste per capita is encoded as hue using a
perceptually-ordered blue-to-red scale.
Choropleth was chosen over proportional symbols because the quantity is an
areal density (kg/person/year), not an absolute count; encoding an areal
variable as point size conflates national area with the variable of interest.
The colour scale is clipped at the 95th percentile to prevent the most
extreme outliers from washing out the majority of countries.

\textbf{Ranked bar chart.}
A horizontal bar chart of the top-$N$ countries (user-controlled, 10--50)
complements the map.
Whereas the choropleth reveals broad geographic patterns at a glance, precise
rank ordering is cognitively difficult from colour alone; the bar chart
resolves this using the most accurate pre-attentive channel (position on a
common scale) [6].

\textbf{Interaction design.}
A year slider (1990--2019) enables temporal exploration; both map and chart
update reactively.
A top-$N$ slider controls the bar chart length.
A multi-select dropdown (selectize widget) allows filtering both views to a
custom country set.
Multi-select was preferred over free-text search because it provides
autocomplete from the full country list, prevents typos, keeps the active
filter set explicit, and supports comparison of non-adjacent countries
(e.g.\ Philippines vs.\ Germany).

\subsection{Act II --- River Emission Pathways}

\textbf{Scatter-geo map.}
Each of the 32\,000 river records is rendered as a circle on a geographic
map.
Marker radius is proportional to $\log_{e}(1 + \text{emissions})$.
Logarithmic scaling is essential: raw emissions span four orders of magnitude,
and a linear encoding would render more than 90\,\% of rivers as
invisibly small dots while a handful of extreme rivers dominate the visual field.

\textbf{K-Means cluster colouring.}
The ten K-Means clusters are each assigned a distinct categorical colour.
Colour here encodes cluster membership (nominal), not emission magnitude
(quantitative).
This choice reflects the primary message of this act: geographic
\emph{hotspot structure}, not relative magnitude.
Magnitude is communicated by the sidebar percentage readouts and the
by-country and by-continent bar charts beneath the map, which provide the
quantitative precision that colour cannot.

\textbf{Interaction design.}
A ``Rivers to show'' slider controls how many of the highest-emitting rivers
are rendered, allowing users to focus on the most impactful sources.
A minimum-emissions threshold slider filters out noise from low-emitting
rivers.
The cluster legend in the sidebar lists each cluster's dominant country,
total percentage of global emissions, and top-3 contributing nations,
making every colour immediately interpretable without hovering.

\subsection{Act III --- Ocean Health Explorer}

\textbf{Choropleth map.}
OHI index scores (0--100) are mapped with a red-to-green diverging scale,
exploiting the culturally universal red\,=\,danger, green\,=\,good
semantic [6].
A year slider (2012--2023) and a goal selector (composite Index or any of
10 sub-goals) allow targeted exploration of specific health dimensions over
time.

\textbf{Trend line chart.}
A line chart for the selected region and goal shows annual OHI scores from
2012 to 2023.
The selected year is highlighted with a coloured diamond marker, linking the
chart to the map's current state.
Fill under the line (with low opacity) provides a quick impression of the
score level without visual overload.

\textbf{Radar chart.}
A ten-axis spider chart shows all OHI sub-goal scores simultaneously for the
selected region and year.
The radar form was chosen over a small-multiples bar chart because the
enclosed area and polygon shape give an immediate holistic fingerprint of a
region's multi-dimensional health profile.
Outlier sub-goals (very low or very high) are immediately visible as
concavities or protrusions in the polygon.

\textbf{Forecast panel.}
A fourth chart extends the trend line with a user-selected model and horizon.
Confidence bands (80\,\% coverage) are shaded to make predictive uncertainty
explicit and to discourage naive extrapolation.

\textbf{Interaction design.}
A region/country dropdown (220+ entries, searchable) links all three charts.
A forecast model selector and horizon slider give users agency over the
predictive component.
In-chart annotation of $R^2$ and MAE supports informed model comparison.

\subsection{Visual Design Rationale}

A dark ocean theme (primary background \texttt{\#080d1a}) was adopted for two
mutually reinforcing reasons: the dark ocean fill provides perceptual
continuity between the map's ocean areas and the application's chrome, and
luminous categorical colours (cyan \texttt{\#00d4ff}, coral
\texttt{\#ff6b35}, green \texttt{\#00ff88}) achieve higher contrast and
distinctiveness against dark backgrounds than against white.
Glassmorphism card borders (\texttt{backdrop-filter: blur}) visually separate
chart panels while maintaining depth transparency.
All three acts use an identical layout template (sticky sidebar + scrollable
main area) to reduce the user's cognitive re-orientation cost when switching
acts.

%% ────────────────────────────────────────────────────────────────────────────
\section{Machine Learning Methods}

\subsection{Emission-Weighted K-Means Clustering}

\textbf{Motivation.}
Displaying all 32\,000 river points simultaneously creates severe visual
clutter and obscures geographic structure.
Clustering reduces this to ten interpretable groups, each summarised by a
centroid, emission percentage, and dominant nation.

\textbf{Method.}
\tool{sklearn.cluster.KMeans} with $k = 10$ is applied to the
(latitude, longitude) coordinate pairs.
Sample weights are set to:
\[
  w_i = \sqrt{\,\max(e_i,\;10^{-6})\,}
\]
where $e_i$ is the river's annual emission in tonnes.
Square-root weighting ensures that high-emitting rivers pull centroids toward
true hotspot cores without allowing a small number of extreme values to
collapse all ten clusters into a single high-emission region.
Linear weighting was tested but placed nine of ten centroids within
South-East Asia; the square-root parameterisation distributes them globally
while still reflecting emission importance.

\textbf{Hyperparameter selection.}
The choice of $k = 10$ was guided by an elbow analysis on within-cluster sum
of squares (WCSS), which showed diminishing returns after $k \approx 8$.
Ten was selected as it aligns with broad geographic regions (South-East Asia,
South Asia, East Asia, West Africa, Latin America, etc.) and fits comfortably
in a sidebar legend.
The random seed is fixed at 42 for reproducibility.

\textbf{Cluster summaries.}
For each cluster the system reports: emission-weighted centroid coordinates,
total annual emissions and percentage of global river load, dominant country,
continent, and the top-3 contributing nations with individual emission totals.
This makes the clusters analytically useful, not merely decorative.

\subsection{Time-Series Forecasting of OHI Scores}

Four models provide point forecasts and 80\,\% confidence intervals for
user-specified horizons of 1--10 years.
A minimum of five historical observations is required to prevent degenerate
fits.
Table~\ref{tab:models} summarises each model and its key assumptions.

\begin{table}[ht]
\centering
\small
\caption{Forecast models, assumptions, and CI method.}
\label{tab:models}
\resizebox{\columnwidth}{!}{%
\begin{tabular}{@{}p{2.3cm}p{2.5cm}p{1.8cm}@{}}
\toprule
\textbf{Model} & \textbf{Assumption} & \textbf{CI source}\\
\midrule
Linear regression & Constant linear trend & Residual std\\
Polynomial (deg.~2) & Quadratic trend & Residual std\\
Holt-Winters (ETS) & Damped additive trend, no seasonality & Residual std\\
ARIMA(1,1,1) & AR(1) on first-differenced series & Analytic 80\%\\
\bottomrule
\end{tabular}}
\end{table}

\textbf{Model justifications.}
\emph{Linear regression} provides a transparent baseline.
\emph{Polynomial regression} (degree~2) captures U-shaped or parabolic
trajectories observed in some recovering/declining regions.
\emph{Holt-Winters with damped trend} (Statsmodels
\tool{ExponentialSmoothing}) handles the near-stationary behaviour typical
of OHI scores: the damping parameter shrinks the projected trend toward the
long-run mean, producing conservative forecasts appropriate for slowly
evolving ecosystem health metrics.
\emph{ARIMA(1,1,1)} (Statsmodels \tool{ARIMA}) addresses mild
non-stationarity via first differencing (order $d = 1$), with one
autoregressive and one moving-average term to model short-term momentum [7].
Providing all four models side-by-side is intentional: it encourages
critical thinking about how different assumptions produce divergent futures
rather than presenting a single ``correct'' prediction.

\textbf{Metrics.}
In-sample $R^2$ and MAE are computed for every model run and displayed in
the chart annotation, enabling direct model comparison without navigating
away from the visualisation.

%% ────────────────────────────────────────────────────────────────────────────
\section{Findings and Insights}

\subsection{Act I --- Plastic Mismanagement}

The choropleth map immediately reveals that mismanagement is a
\emph{geographically concentrated} problem.
South and South-East Asian countries --- the Philippines, India, Algeria,
and Vietnam --- consistently occupy the top five positions across all
available years.
This reflects rapid urbanisation combined with waste-collection
infrastructure that has not scaled proportionally.
High-income nations (Western Europe, North America, Japan) register near
0--1\,kg/capita/year throughout, underscoring that the crisis is driven
by development gaps rather than total production volume alone.

Temporal exploration with the year slider reveals a \emph{plateau effect}
in many middle-income nations post-2010, consistent with international
waste-management investment during this period.
However, the most severely affected nations show no substantial improvement
over the 30-year window, suggesting that infrastructure investment has not
reached the communities with the highest mismanagement rates.

\subsection{Act II --- River Emission Pathways}

The scatter-geo map confirms the extreme Pareto concentration reported by
Meijer et al.\ [2]: the top 1\,\% of rivers by emission volume contribute
over 60\,\% of total global riverine plastic loading.
Seven of the ten K-Means clusters are centred on Asian river systems
(Mekong, Ganges, Yangtze, Pasig, and their tributaries), with the remaining
two clusters on West African and Latin American rivers respectively.

Cross-referencing Acts~I and~II shows strong geographic alignment:
the Philippines, India, and China dominate both the waste mismanagement
rankings and the river loading rankings.
This co-occurrence, while not a statistical proof of causation, is
consistent with the hypothesis that poor land-based waste management
is the proximate driver of river loading.
It also highlights the policy leverage point: interventions that reduce
mismanaged waste at source (waste collection, recycling infrastructure)
should produce downstream reductions in river and ocean contamination.

\subsection{Act III --- Ocean Health Trends and Forecasts}

The OHI choropleth reveals that aggregate ocean health is neither uniformly
good nor catastrophically bad: the global average composite Index score
fluctuates in the narrow band 69--72 across the entire 2012--2023 measurement
window.
This apparent stability, however, conceals significant \emph{sub-goal
heterogeneity} visible in the radar chart.
Coastal Protection and Carbon Storage regularly score above 80 for mid-ocean
island states, while Livelihoods \& Economies and Artisanal Opportunities
frequently fall below 60 in densely populated coastal nations in South Asia
and West Africa.

Regional divergence is stark: small Pacific island states (Cocos Islands,
Palmyra Atoll) consistently achieve scores above 85, reflecting minimal
anthropogenic pressure, while regions adjacent to major Asian river mouths
(Bay of Bengal, South China Sea) show below-average and in some sub-goals
declining scores, plausibly linked to the riverine plastic loading revealed
in Act~II.

All four forecast models agree on continued modest stability or a slight
downward trajectory for most regions through 2028.
ARIMA(1,1,1) produces the widest confidence bands, reflecting its
sensitivity to recent one-year fluctuations.
Holt-Winters consistently produces the most conservative (lowest slope)
extrapolations due to trend damping, which users often find most credible
for near-stationary ecological metrics.

%% ────────────────────────────────────────────────────────────────────────────
\section{Challenges}

\textbf{OHI region--country mismatch.}
The OHI dataset defines 220+ ocean regions (e.g.\ ``Northeast US'',
``California Current'', ``Bay of Bengal'') that correspond to oceanographic
areas, not to ISO country boundaries.
The React/D3 frontend resolves this via the official OHI GeoJSON boundary
file, which provides polygon geometries keyed by \tool{region\_id}.
The Python Shiny frontend falls back to Plotly's
\tool{locationmode='country names'}, which successfully renders approximately
70\,\% of regions (those whose names match standard country names) but leaves
the remaining $\approx$30\,\% (sub-national and ocean-zone regions) unmapped.
A complete fix requires passing the OHI GeoJSON as a
\tool{geojson=} argument to \tool{px.choropleth}, keyed on \tool{region\_id}.

\textbf{Geospatial dependency chain.}
The river pipeline requires GDAL, Fiona, PyProj, and geopandas ---
a dependency footprint of $\approx$300\,MB.
Cold-start time (first spatial join + K-Means clustering) is approximately
30 seconds.
Caching the processed output as a \tool{.parquet} file at first run reduces
subsequent starts to under 2 seconds.
This optimisation was deferred to keep the codebase self-contained and
reproducible.

\textbf{Shiny layout constraints.}
The default \tool{fillable=True} mode in bslib's \tool{layout\_sidebar}
distributes all child elements equally across the viewport height.
With three to four stacked chart sections, each receives only 25\,\% of
the screen, making all charts illegibly small.
No first-class ``scrollable main panel'' API existed at the time of
development; the solution required custom CSS overrides:
\tool{position:sticky} on the sidebar, \tool{overflow-y:auto} on the main
panel, and \tool{flex-shrink:0} on all card elements.

\textbf{Temporal data gap.}
The plastic waste dataset ends in 2019; OHI scores begin in 2012.
The five-year overlap (2014--2019) is too short for statistically robust
cross-dataset correlation analysis between waste mismanagement rates and
OHI score changes.

%% ────────────────────────────────────────────────────────────────────────────
\section{Future Work}

\textbf{Cross-act linked selection.}
The most impactful enhancement would be globally linking selection state:
clicking a country in Act~I should highlight its rivers in Act~II and
navigate to its OHI region in Act~III.
This would transform three adjacent visualizations into a single
integrated analytical system.

\textbf{Complete OHI choropleth.}
Replacing the country-name lookup with the official OHI GeoJSON boundary
file would render all 220+ ocean regions correctly, including sub-national
coastal zones and open-ocean regions.

\textbf{Statistical correlation.}
Computing Pearson and Spearman correlations between a country's per-capita
plastic waste and the OHI score of its adjacent ocean region, with bootstrap
confidence intervals, would test the implied causal link narratively
established by the three-act structure.

\textbf{Ocean current drift.}
Overlaying oceanographic drift trajectory data would visually connect
river outflow points to the downstream ocean zones they pollute, providing
the missing mechanistic link between Acts~II and~III.

\textbf{Prophet forecasting.}
Re-enabling the Prophet model [currently excluded due to the \tool{cmdstanpy}
compilation requirement] would add automatic change-point detection to the
forecast suite, potentially capturing regime shifts in OHI trends.

%% ────────────────────────────────────────────────────────────────────────────
\section{Conclusion}

This project demonstrates that a structured three-act narrative, combined with
emission-weighted K-Means clustering and a four-model forecast suite, enables
users across expertise levels to trace the full pathway of plastic pollution ---
from land-based mismanagement, through river transport, to ocean health
consequences.
Key findings confirm that plastic mismanagement is geographically concentrated
in South and South-East Asia, that river loading follows an extreme Pareto
distribution with the top 1\,\% of rivers carrying the majority of the load,
and that ocean health has remained in a narrow stability band over the 12-year
measurement window despite growing contamination.
The dual-implementation strategy (React/D3 for production deployment and
Python Shiny for research prototyping) validates that the same data pipeline,
preprocessing logic, and ML methods can serve both deployment contexts with
minimal code duplication.

%% ────────────────────────────────────────────────────────────────────────────
{\small
\begin{thebibliography}{9}

\bibitem{jambeck2015}
J.\,R. Jambeck et~al.,
``Plastic waste inputs from land into the ocean,''
\textit{Science}, vol.~347, no.~6223, pp.~768--771, 2015.

\bibitem{meijer2021}
L.\,J.\,J. Meijer et~al.,
``More than 1000 rivers account for 80\% of global riverine plastic emissions
into the ocean,''
\textit{Science Advances}, vol.~7, no.~18, p.~eaaz5803, 2021.

\bibitem{halpern2012}
B.\,S. Halpern et~al.,
``An index to assess the health and benefits of the global ocean,''
\textit{Nature}, vol.~488, pp.~615--620, 2012.

\bibitem{ritchie2018}
H. Ritchie and M. Roser,
``Plastic pollution,''
\textit{Our World in Data}, 2018.
\url{https://ourworldindata.org/plastic-pollution}

\bibitem{segel2010}
E. Segel and J. Heer,
``Narrative visualization: Telling stories with data,''
\textit{IEEE Transactions on Visualization and Computer Graphics},
vol.~16, no.~6, pp.~1139--1148, 2010.

\bibitem{munzner2014}
T. Munzner,
\textit{Visualization Analysis and Design}.
AK Peters/CRC Press, 2014.

\bibitem{box2015}
G.\,E.\,P. Box, G.\,M. Jenkins, G.\,C. Reinsel, and G.\,M. Ljung,
\textit{Time Series Analysis: Forecasting and Control}, 5th~ed.
Wiley, 2015.

\end{thebibliography}
}

\end{document}
